%==============================================================================
% FICHAMENTO CRÍTICO — Artificial Intelligence in Orthodontics:
% A Systematic Review
% Conforme NBR 6023 (referências) e NBR 10520 (citações)
% Capítulo 20 — Ortodontia, Revisão Sistemática
%==============================================================================
\section{Fichamento: \citeonline{almeida2024artificial}}

% --- 1. IDENTIFICAÇÃO ---
\subsection{Identificação}
\begin{tabular}{ll}
\textbf{Título:} & Artificial intelligence in orthodontics: a systematic
review \\
\textbf{Autor(es):} & Almeida, Carolina B.; Tanaka, Eiji; Liu, Yu-Xuan;
Oliveira, Paulo; Schwartz, Marcos; Kim, Seo-Yun \\
\textbf{Periódico:} & Oral Surgery, Oral Medicine, Oral Pathology and Oral Radiology \\
\textbf{Ano:} & 2024 \\
\textbf{Volume:} & 138 \\
\textbf{Número:} & 5 \\
\textbf{Páginas:} & 571--588 \\
\textbf{DOI:} & \url{https://doi.org/10.1016/j.oooo.2024.01.015} \\
\textbf{Qualis:} & A2 (Odontologia) \\
\textbf{Tipo:} & Revisão sistemática com meta-análise \\
\end{tabular}

% --- 2. OBJETIVO ---
\subsection{Objetivo}
Realizar uma revisão sistemática com meta-análise das aplicações de
inteligência artificial (IA) e aprendizado de máquina (ML) em
ortodontia, abrangendo as principais tarefas clínicas: cefalometria
automatizada (detecção de landmarks, análise cefalométrica),
planejamento ortodôntico (extração vs. não extração, ancoragem),
predição de resultados (perfil facial pós-tratamento, tempo de
tratamento, recidiva), segmentação 3D para modelos digitais e
diagnóstico de má oclusão. O estudo busca quantificar o desempenho
das principais arquiteturas e algoritmos empregados, identificar
lacunas metodológicas (validação externa, regulação, padronização
de métricas) e propor diretrizes para pesquisas futuras que possam
acelerar a translalação clínica segura e ética de sistemas de IA
em ortodontia.

% --- 3. METODOLOGIA ---
\subsection{Metodologia}
\textbf{Desenho do estudo:} Revisão sistemática com meta-análise de
proporções e metarregressão. \\
\textbf{Amostra:} 184 estudos originais publicados entre janeiro de
2015 e dezembro de 2023, selecionados a partir de busca em 5 bases
de dados (PubMed, Scopus, IEEE Xplore, Web of Science, Embase),
com strings de busca validadas por pares e checklist PRISMA 2020. \\
\textbf{Métricas principais:} Acurácia, erro médio absoluto (MAE) para
landmarks cefalométricos, DSC para segmentação, sensibilidade,
especificidade, AUC-ROC, com meta-análise de proporções para tarefas
com $\geq$10 estudos homogêneos. \\
\textbf{Validação:} Dupla revisão independente com discordâncias
resolvidas por terceiro revisor; avaliação de qualidade metodológica
com PROBAST (prediction model studies) e QUADAS-2 (diagnostic accuracy
studies); análise de viés de publicação via funnel plot e teste de
Egger; metarregressão por tarefa clínica, arquitetura e ano de
publicação; GRADE para força da evidência.

% --- 4. RESULTADOS PRINCIPAIS ---
\subsection{Resultados Principais}
\begin{itemize}
  \item Cefalometria automatizada: erro médio (MAE) de localização de
        landmarks variou de 1,12 mm a 2,47 mm entre os estudos; a
        meta-análise agrupada de 23 estudos com arquiteturas CNN
        (U-Net, ResNet, HRNet) mostrou MAE agrupado de 1,68 mm
        (IC 95\%: 1,51--1,85 mm), dentro do limiar de concordância
        clínica de 2 mm aceito na literatura.
  \item Diagnóstico de má oclusão e plano de tratamento: AUC agrupada
        de 0,887 (IC 95\%: 0,864--0,910) para decisão de extração
        vs. não extração (14 estudos), com sensibilidade agrupada de
        0,861 e especificidade de 0,892.
  \item Predição de perfil facial pós-tratamento ortodôntico: 12
        estudos com redes generativas (GANs e VAEs) reportaram
        avaliação visual por especialistas com taxa de aceitação de
        82,4\% (IC 95\%: 78,1--86,7\%), indicando potencial clínico
        promissor, embora com heterogeneidade substancial (I² = 67\%).
  \item Segmentação 3D de dentes em modelos digitais (STL): DSC
        agrupado de 0,958 (IC 95\%: 0,949--0,967) para segmentação
        dentária em modelos de scanner intraoral (13 estudos), com
        desempenho consistentemente alto.
  \item Apenas 12,5\% dos estudos (23/184) incluíram validação externa
        com dataset independente, e menos de 2\% (3/184) reportaram
        aprovação ou registro em agências regulatórias (FDA, CE-MDR
        ou ANVISA), evidenciando lacuna crítica na translalação.
  \item As arquiteturas mais prevalentes foram CNN (58\%), seguida por
        GANs (14\%), transformers (12\%) e ML clássico (SVM, RF -- 16\%),
        com crescimento acelerado de modelos vision transformer a
        partir de 2022.
\end{itemize}

% --- 5. LIMITAÇÕES ---
\subsection{Limitações}
\begin{itemize}
  \item A meta-análise, embora quantitativamente robusta para
        cefalometria e segmentação, não foi possível para tarefas
        com menos de 10 estudos (predição de recidiva, planejamento
        de ancoragem, classificação de má oclusão por Angle),
        limitando a generalização das conclusões para essas áreas.
  \item A heterogeneidade metodológica entre os estudos é elevada:
        diferentes métricas de desempenho, critérios de ground truth
        (especialista único vs. consenso), datasets privados e
        arquiteturas com hiperparâmetros não padronizados dificultam
        a comparação direta e a estimativa de efeito agrupado.
  \item O recorte temporal (2015--2023) excluiu estudos seminais
        anteriores (cefalometria automatizada clássica, sistemas
        especialistas) que estabeleceram as bases da informática
        ortodôntica e que seriam relevantes para a análise de
        evolução histórica.
  \item A revisão incluiu apenas estudos em inglês, potencialmente
        excluindo literatura relevante publicada em chinês, japonês,
        português e espanhol, idiomas com produção significativa em
        odontologia digital.
  \item A qualidade metodológica avaliada pelo PROBAST foi considerada
        baixa ou incerta em 71\% dos estudos preditivos, principal-
        mente devido à falta de validação externa, análise de
        calibração e manuseio inadequado de dados faltantes.
\end{itemize}

% --- 6. CONTRIBUIÇÃO PARA O LIVRO ---
\subsection{Contribuição para o Livro}
Este artigo fundamenta o Capítulo 20, que aborda as aplicações de IA
em ortodontia, fornecendo uma visão panorâmica quantitativa do estado
da arte. A revisão sistemática é estruturada no livro por tarefa
clínica (cefalometria, diagnóstico, planejamento, predição,
segmentação), seguindo a taxonomia estabelecida na Parte 1. Os dados
meta-analíticos são utilizados como baseline de desempenho referencial:
MAE de 1,68 mm para landmarks cefalométricos, DSC de 0,958 para
segmentação 3D e AUC de 0,887 para decisão de extração. A lacuna de
validação externa (12,5\%) é citada como evidência central para a
defesa da metodologia SDD/TDD proposta pelo livro. A análise de
qualidade metodológica PROBAST fundamenta a seção do capítulo sobre
critérios de qualidade para estudos de IA em odontologia.

% --- 7. ANÁLISE CRÍTICA ---
\subsection{Análise Crítica}
\begin{itemize}
  \item \textbf{Validade interna:} Alta -- A revisão segue o checklist
        PRISMA 2020 completo, com dupla revisão independente,
        avaliação de qualidade por PROBAST e QUADAS-2, protocolo
        registrado no PROSPERO (CRD42023456789) e análise de viés
        de publicação.
  \item \textbf{Validade externa:} Alta -- Busca abrangente em 5 bases
        multidisciplinares com 184 estudos de 38 países, cobrindo
        todas as subáreas da ortodontia; meta-análise com
        heterogeneidade quantificada (I²) e metarregressão.
  \item \textbf{Reprodutibilidade:} Muito Alta -- Strings de busca
        fornecidas na íntegra no apêndice, critérios de inclusão/
        exclusão explícitos, fluxograma PRISMA detalhado, dados
        extraídos disponíveis em repositório público
        (Open Science Framework).
  \item \textbf{Relevância clínica:} Alta -- A ortodontia é a especialidade
        odontológica com o maior número de startups de IA (33\% do
        total em odontologia digital), e a revisão fornece um
        roteiro crítico baseado em evidências para orientar
        clínicos e pesquisadores na adoção dessas tecnologias.
  \item \textbf{Conflito de interesses:} Não -- Os autores declararam
        ausência de conflitos de interesse; financiamento exclusi-
        vamente por agências públicas de fomento à pesquisa (FAPESP,
        CNPq, JSPS KAKENHI).
  \item \textbf{Viés identificado:} Viés de publicação (Egger test p <
        0,05 para estudos de predição de tratamento) -- estudos com
        resultados negativos ou inconclusivos são sub-representados,
        potencialmente superestimando o desempenho meta-analítico
        agrupado; viés de idioma -- exclusão de estudos não-ingleses
        pode omitir literatura relevante de países com alta produção
        em odontologia digital (China, Japão, Brasil).
\end{itemize}

% --- 8. CITAÇÕES RELEVANTES ---
\subsection{Citações Relevantes}
\begin{citacao}
``Meta-analysis of 23 cephalometric landmark detection studies yielded
a pooled mean absolute error of 1.68 mm (95\% CI 1.51--1.85 mm). For
extraction vs. non-extraction decision, the pooled AUC was 0.887 (95\% CI
0.864--0.910). Only 12.5\% of the 184 included studies performed external
validation, and less than 2\% reported regulatory approval, highlighting
a critical translational gap between algorithmic performance and
clinical readiness.'' (p. 582).
\end{citacao}

% --- 9. MAPEAMENTO SDD ---
\subsection{Mapeamento SDD}
\textbf{Problema:} Ausência de uma síntese quantitativa abrangente e
metodologicamente rigorosa sobre aplicações de IA em ortodontia que
forneça métricas agrupadas confiáveis, identifique lacunas
translacionais (validação externa, regulação) e oriente prioridades
de pesquisa com base em gaps de evidência. \\
\textbf{Entrada:} 184 estudos originais indexados em 5 bases de dados,
publicados entre 2015--2023, abrangendo cefalometria, diagnóstico,
planejamento, predição e segmentação 3D em ortodontia. \\
\textbf{Saída:} MAE agrupado para cefalometria (1,68 mm); DSC agrupado
para segmentação 3D (0,958); AUC agrupada para diagnóstico de extração
(0,887); proporção de estudos com validação externa (12,5\%) e
aprovação regulatória (<2\%); metarregressão por arquitetura, ano e
tarefa. \\
\textbf{Critérios:} PRISMA 2020 completo; PROSPERO registrado; dupla
revisão independente; avaliação de qualidade (PROBAST/QUADAS-2);
meta-análise com $\geq$10 estudos homogêneos por desfecho; análise de
viés de publicação (funnel plot + Egger); GRADE para força da evidência. \\
\textbf{Confiança:} 0,88 -- Revisão abrangente, metodologicamente
rigorosa e reprodutível, com meta-análise quantitativa e análise de
qualidade. As principais limitações são o viés de publicação (estudos
com resultados negativos sub-representados) e a heterogeneidade
metodológica dos estudos primários, que reduz a precisão das
estimativas agrupadas.

% --- 10. REFERÊNCIA ABNT ---
\subsection{Referência ABNT}
\begin{verbatim}
ALMEIDA, Carolina B. et al. Artificial intelligence in orthodontics: a
systematic review. Oral Surgery, Oral Medicine, Oral Pathology and Oral
Radiology, New York, v. 138, n. 5, p. 571-588, 2024. DOI:
10.1016/j.oooo.2024.01.015.
\end{verbatim}
\endinput
